% Part: proof-theory
% Chapter: proof-search
% Section: introduction

\documentclass[../../../include/open-logic-section]{subfiles}

\begin{document}

\olfileid{pt}{ps}{int}

\olsection{پېژندنه}

د \Log{G3c} يا \Log{N1i} په شان د !!{proof} نظامونو
بديهيتونه او استنتاجي قواعد ټاکي چې د سيکوېنټونو يا
!!{formula}s کومې ونې صحيح !!{proof}s ګڼل کېږي۔ که د
سيکوېنټونو يا !!{formula}s داسې يوه ونه راکړل شي، کېدے شي
ورسره نور معلومات هم وي: په هر ګام کښې کوم قواعد کارېدلي،
يا د فرضونو او هغو استنتاجي قواعدو نښې چې دغه فرضونه باسي۔
د بديهيتونو او قواعدو تعريفونه بيا دا کتنه ممکنوي چې
راکړې ونه صحيح !!{proof} ده که نۀ۔ خو د يو راکړي سيکوېنټ،
يا له راکړو فرضونو څخه د يو راکړي !!{formula} دپاره
!!a{proof} \emph{موندل} بېله او عموماً ګرانه مسئله ده۔
همدغه د \emph{ثبوت د لټون} مسئله ده۔

د ثبوت د لټون يوه ډېره ساده طريقه شته۔ !!{formula}s د
متناهي الفبا د نښو لړۍ ګڼلے شو۔ سيکوېنټونه او د
سيکوېنټونو يا !!{formula}s ونې هم د !!{formula}s لړۍ
ګڼلے شو؛ کېدے شي د ونې د څانګېدو دپاره ځانګړي قوسونه
وکاروو۔ بديهيتونه او قواعد اجازه راکوي چې په ميخانيکي
ډول وګورو آيا د نښو يوه لړۍ صحيح !!{proof} ښيي۔ نو په
ميخانيکي ډول \emph{ټول ممکن} صحيح !!{proof}s جوړولے شو:
د !!{proof}s د ښودلو د الفبا ټولې نښيزې لړۍ جوړوو او د هرې
ممکنې لړۍ دپاره ګورو چې صحيح !!{proof} ښيي که نۀ۔ د
ثبوت د لټون ساده طريقه دا ده: ټول صحيح !!{proof}s تر
هغه وخته جوړ کړئ چې له هغو يو د راکړي سيکوېنټ، يا د
راکړي !!{formula}، !!{proof} وي او پرانيستي فرضونه ئې په
راکړو فرضونو کښې شامل وي۔

غوره طريقه هغه ساده لاره ده چې په لاس د !!a{proof} د موندلو
په وخت ئې کاروئ: له پاييز سيکوېنټ څخه پيل کوئ، يو
!!a{formula} ټاکئ او استنتاجي قاعده ``شاته'' کاروئ، څو
يوه يا څو مقدمې جوړې شي۔ که د دغو مقدمو !!{proof}s موجود
وي، له وروستي استنتاج سره يوځاے د راکړي سيکوېنټ
!!a{proof} جوړوي۔ په طبيعي استنتاجي نظام کښې هم د
!!a{proof} د موندلو دپاره ورته، خو يو څه پېچلې، طريقه
کاروئ۔

خو دغه طريقه بې‌خطا نۀ ده: که ناسمه قاعده وټاکئ، يا
!!{formula}s په ناسم ترتيب وټاکئ، بې‌لارې پاے ته رسېدې
شئ۔ دا طريقه پوره ټاکلې هم نۀ ده: په هر ځاے کښې کېدے
شي د !!{formula} د ټاکلو يا د قاعدې د شاته کارولو څو
امکانونه وي۔ د \emph{ثبوت د لټون الګوريتم} يوه پوره
ټاکلې کړنلاره ده چې که !!a{proof} موجود وي، هغه مومي؛
يعنې بې‌خطا ده۔

د ځينو !!{proof} نظامونو دپاره د ثبوت د لټون ساده
الګوريتمونه جوړېږي۔ په سيکوېنټي نظامونو کښې د
!!a{formula} !!{main operator} او دا چې په کيڼ يا ښي اړخ
کښې دے، ټاکي چې کومه قاعده شاته وکارول شي۔ د بېلګې په
توګه، که د $\Gamma \Sequent \Delta, !A \land !B$
!!a{proof} غواړئ، چې $!A \land !B$ ئې په ښي اړخ کښې
دے، نو \RightR{\land} شاته کاروئ او دوه اړوندې مقدمې
$\Gamma \Sequent \Delta, !A$ او $\Gamma \Sequent \Delta, !B$
جوړوئ۔ خو که نظام مو \Log{G1c} وي، د راغونډولو موجودګي
کار ګرانوي، ځکه نور امکانونه درکوي: \RightR{\Contraction}
هم شاته کارولے او مقدمه
$\Gamma \Sequent \Delta, !A \land !B, !A \land !B$
جوړولے شئ؛ بيا
$\Gamma \Sequent \Delta, !A \land !B, !A \land !B, !A \land !B$
او همداسې نور۔ په \Log{G2c} کښې کار لا ګران دے۔ د
\RightR{\land} شاته کارونه دا هم غواړي چې $\Gamma_1$،
$\Gamma_2$ داسې اټکل کړئ چې $\Gamma = \Gamma_1 \cup \Gamma_2$
وي، او $\Delta_1$، $\Delta_2$ داسې چې
$\Delta = \Delta_1 \cup \Delta_2$ وي؛ بيا د مقدمو
$\Gamma_1 \Sequent \Delta_1, !A$ او
$\Gamma_2 \Sequent \Delta_2, !B$ هڅه وکړئ۔ ناسم اټکل
وروسته بې‌لارې پاے ته رسولے شي؛ بيا بايد پيل ته بېرته
لاړ شئ او نور $\Gamma_1$، $\Gamma_2$، $\Delta_1$،
$\Delta_2$ وټاکئ۔ که !!{formula} د $!A \lor !B$ په بڼه
وي، بايد د \RightR{\lor} له دوو امکانونو څخه يو وټاکئ؛
دا يا مقدمه $\Gamma \Sequent \Delta, !A$ او يا مقدمه
$\Gamma \Sequent \Delta, !B$ جوړوي۔ ناسم انتخاب بيا هم
شاته تګ ته اړکولے شي۔

نظام \Log{G3c} دغه ستونزې نۀ لري۔ جوړښتي قواعد نۀ لري،
او د استنتاجي قاعدې انتخاب د !!{formula} په
!!{main operator} او په هغه اړخ ټاکل کېږي چې فارمول
پکښې موجود دے۔ نو \Log{G3c} د ثبوت د لټون دپاره له
\Log{G1c} يا \Log{G2c} څخه زيات مناسب دے۔ بريالے لټون
په \Log{G3c} کښې !!a{proof} ورکوي؛ دغه !!{proof} بيا په
\Log{G1c} يا \Log{G2c}، او حتٰی په \Log{N1c}، کښې
!!{proof}s ته اړولے شو۔ نو که د \Log{G3c} د ثبوت د
لټون الګوريتم او د \Log{G3c} !!{proof}s نورو نظامونو ته
د اړولو الګوريتمونه ولرو، د دغو نورو نظامونو د ثبوت د
لټون طريقې هم ترلاسه کېږي۔

څنګه پوهېږو چې د ثبوت د لټون الګوريتم بې‌خطا دے؟ بايد
وښيو چې که الګوريتم !!a{proof} نۀ مومي، هېڅ !!{proof}
نۀ شي موجودېدې۔ که ناسم الګوريتم وټاکو، کېدے شي داسې
حالتونه راشي چې د موجود ثبوت له موندلو ناکام شي۔ ساده
الګوريتم دغه ستونزه نۀ لري، ځکه د ټولو ممکنو
!!{proof}s لټون کوي۔ خو د ثبوت د لټون الګوريتم بايد
اغېزمن وي او لږ تر لږه هڅه وغواړي۔

د ثبوت د لټون د بې‌خطا والي د ښودلو بنسټيزه مفکوره
دا ده: که الګوريتم ناکام شي، د ناکامۍ له طريقې ښيو چې
سيکوېنټ هېڅ !!a{proof} نۀ شي لرلے۔ دا په دې ښيو چې
سيکوېنټ معتبر نۀ دے۔ په کلاسيکي لومړۍ درجې منطق کښې
معتبر سيکوېنټونه $\Gamma \Sequent \Delta$ هغه دي چې
په هر !!{structure} کښې يا د $\Gamma$ لږ تر لږه يو
!!{formula} دروغ وي، يا د~$\Delta$ لږ تر لږه يو
!!{formula} رښتيا وي۔ \Log{G3c} \emph{سم} دے؛ يعنې
يوازې معتبر سيکوېنټونه !!{prove}s کوي۔ دا د !!{proof}s پر
سر په استقرا ثابتېږي: بديهيتونه ښکاره معتبر دي، او که
د يوې قاعدې مقدمې معتبرې وي، پايله ئې هم معتبره ده۔
د الګوريتم د بې‌خطا والي د ښودلو دپاره ښيو چې که د
$\Gamma \Sequent \Delta$ د !!a{proof} لټون ناکام شي،
!!a{structure} تعريفولے شو چې پکښې د $\Gamma$ ټول
!!{formula}s رښتيا او د~$\Delta$ ټول !!{formula}s دروغ
وي۔ دا د \Log{G3c} \emph{بشپړتيا} هم ثابتوي؛ يعنې که
$\Gamma \Sequent \Delta$ معتبر وي، د هغه يو
\Log{G3c}-!!{proof} موجود وي۔ ځکه چې د ثبوت هر ناکام
لټون د $\Gamma \Sequent \Delta$ ناسم والے ښيي، د معتبر
سيکوېنټ هر لټون بايد بريالے شي۔

\end{document}
